\documentclass[exercice]{thedocument}%
\usepackage{texmaths-path}
\startdatakeys
\source{19GENMATMEAG2}
\cycle{4}
\competencesDuSocle{CA,CH,RA,CO}
\connaissancesRequises{D2f*,D2d1}
\competencesTravaillees{D2e*}
\motsCles{arctangente,Thalès direct}
\enddatakeys
\begin{document}%
Leila est en visite à Paris. Aujourd'hui, elle est au Champ de Mars où l'on
peut voir la tour Eiffel dont la hauteur totale \mathspoint{BH} est 324 m.%
\par\medskip\noindent
Elle pose son appareil photo au sol à une distance $\mathspoint{AB}=600$ m
du monumentet le programme pour prendre une photo (voir le dessin ci-dessous).
\begin{enumerate}
    \item Quelle est la mesure, au degré près, de l'angle $\widehat{\mathspoint{HAB}}$~?
    \item Sachant que Leila mesure 1,70 m, à quelle distance \mathspoint{AL} de son
    appareil doit-elle se placer pour paraître aussi grande que la tour Eiffel sur
    la photo~?\par\noindent
    Donner une valeur approchée du résultat au centimètre près.
\end{enumerate}
{\centering\pic[scale=0.75]{19GENMATMEAG2-ex4-fig1.png}\par}
\showthesource
\end{document}